\begin{figure}[htbp]
\centering
\begin{tikzpicture}[x=5.6mm,y=5.6mm]
  \begin{scope}
    \bcuadricula{4}
    \foreach \r/\c in {1/2,2/4,3/1,4/3} {\breina{4}{\r}{\c}}
    \node[font=\small,anchor=north] at (2,-.35) {(a) Solución};
  \end{scope}
  \begin{scope}[shift={(9,0)}]
    \bcuadricula{4}
    \draw[btinta,densely dashed,line width=1.1pt] (.5,3.5)--(3.5,.5);
    \foreach \r/\c in {1/1,2/2,3/3,4/4} {\breina{4}{\r}{\c}}
    \node[font=\small,anchor=north] at (2,-.35) {(b) Diagonal compartida};
  \end{scope}
\end{tikzpicture}
\caption{Dos permutaciones en el mismo tablero. Las filas se numeran de arriba abajo y las columnas
de izquierda a derecha, empezando en uno. En (a) se leen las columnas $(2,4,1,3)$: las cuatro
reinas están también en diagonales distintas de cada orientación. En (b) se leen $(1,2,3,4)$:
cada fila y cada columna sigue teniendo una reina, pero la diagonal trazada contiene las cuatro.}
\label{fig:permutaciones}
\end{figure}
